\chapter{核心技术要点与难点分析}
\section{多速率信号的统一时间组织}
系统同时处理2 MHz PDM、16 kHz PCM、标称200 Hz XYZ、较慢的温度及异步开关按键。若按电脑收到数据的时间简单排列，会混入串口等待和界面刷新抖动。为此，硬件以100 MHz公共计数器作为统一时基，协议时间单位为10 ns；PC保存板端时刻，并将不同任务与各自数据编号关联。

高频处理使用使能脉冲推进，PCM仍在同一逻辑域处理；XYZ保留实际数据就绪时刻。两路特征以10 ms半开窗口聚合，提供可比较、可回放的数据组织方式。

\fig{fusion}{声振特征、候选配对与前后触发窗口}{图中脉冲和时刻为原理示意。联合触发要求候选间隔不超过设定窗，并进行去重。}

声音特征先以Q8状态估计直流分量，再在含160个PCM点的完整窗口内求去直流幅值的均值：
\begin{equation}
 A_k=\frac{1}{160}\sum_{n\in\mathcal W_k}|x_n-d_n|.
\end{equation}
振动特征采用相邻XYZ的$L_1$差分，并取窗口内最大值：
\begin{equation}
 V_k=\max_{i\in\mathcal W_k}\bigl(|x_i-x_{i-1}|+|y_i-y_{i-1}|+|z_i-z_{i-1}|\bigr).
\end{equation}
触发检测关注阈值上穿，避免持续高电平反复生成同一个候选。联合策略在时间窗内配对两类候选，再用去重窗口抑制同一次活动的重复保存。它属于可解释的决策级融合：每次决定都能追溯到特征、阈值、候选时刻和参数版本。

这一设计的难点是“统一语义”而非简单相加。声音可能来自空气传播，振动可能来自桌面传导；要求两者同时满足能够抑制部分干扰，也可能漏掉只在一路明显的轻敲。前期对照结果用于呈现这种权衡，不能将联合条件天然解释为更高准确率。

\section{有限BRAM下的长录音与抗混叠}
Artix-7内部存储需要同时容纳事件数据、音乐数据、监测快照和协议状态。直接将20秒音乐按16 kHz保存，需要640 kB，仅声音就显著增加资源压力。最终方案保持原16 kHz采集链路，通过31抽头FIR低通后每两点保留一点，生成8 kHz音乐路径。

滤波器使用Hamming窗设计，标称截止频率3.2 kHz，Q15系数和为32768。对称系数对应15个原速样本的群延迟，即0.9375 ms。实现利用100 MHz逻辑时钟与低采样率间的空闲周期串行复用乘加资源，避免为每个抽头铺设独立乘法器。定点累加后做舍入和饱和，输入索引随输出保存，用于核对连续性。

\begin{table}[H]\centering\caption{音乐降采样设计参数（系数计算结果）}\small
\begin{tabularx}{\linewidth}{p{49mm}Y}\toprule
项目 & 参数或结果\\\midrule
输入/输出采样率 & 16 kHz / 8 kHz\\
FIR长度与系数量化 & 31抽头，Q15，系数和32768\\
标称截止频率 & 3200 Hz\\
1 kHz处响应 & 约+0.002 dB\\
6 kHz处响应 & 约$-61.15$ dB\\
线性相位群延迟 & 15个输入样本，0.9375 ms\\
20秒缓冲量 & 160000个PCM16点，共320000 B\\\bottomrule
\end{tabularx}\end{table}

这些响应值来自系数计算，不是模拟仪器测得的板级幅频曲线。抗混叠滤波在抽取前抑制超过新奈奎斯特频率的成分；相较于直接隔点丢样，它有明确的频域依据。原速事件与FFT路径继续保留，长录音优化只作用于音乐数据产品。

\section{可靠传输与可追溯的数据对象}
UART帧包含同步字、协议版本、命令、序号、事件号、长度与CRC32。所有多字节整数采用小端，单帧负载不超过512 B。接收端在字节流中寻找同步字，检查长度和CRC；响应回送序号，使请求与返回可对应。完整记录再使用总长度与总CRC检查，并以SHA-256描述归档文件身份。

\fig{protocol}{UART帧结构及三层核对机制}{CRC覆盖version至payload，不包含同步字。单帧校验、完整记录校验和归档哈希解决不同范围的问题。}

系统保留冻结记录直到主机验证完成，按“事件号+偏移”读取，使重传不会依赖隐含的读指针状态。ACK不仅表示收到响应，而是携带当前记录身份、长度与内容校验后释放存储。断线重读、旧任务中止及迟到响应必须核对编号，从协议层减少跨任务串扰。

工程保留构建编号、固件哈希、配置版本及原始输入，使图表和识别结果可追溯；也便于未来AI评价时区分模型变化与采集链路变化。

\section{三敲唤醒、门控与取消一致性}
敲击检测以XYZ变化量为基础，设置高低阈值形成重新武装条件，并用不应期隔离同一次机械振动。最终录音入口由独立三次候选收集器控制：相邻间隔均符合下限与上限才生成启动脉冲；错误间隔从当前候选重新开始。进入就绪不立即录音，忙碌、切换静默或配置变化会清空不适用的待收集状态。

三敲收集器专门承担启动条件，就绪显示\code{tAP3}并以三次短脉冲提示。拨动开关的瞬态由切换静默门控处理，不改写采集样本。

取消比启动更复杂：录音可能还在板上进行，传输可能进行到一半，识别可能已交给独立进程。系统把取消作为当前任务的状态变更，先处理板端编号，再结束工作进程与旧任务输出。即使新任务已开始，旧任务的迟到结果也必须被身份检查拦截。这种一致性要求在交互平台中与采集算法同样重要。

\section{智能音频识别：特征检索与服务协同}
本地音频指纹以频谱地标为基础，从时频峰值及其相对关系构造哈希，在指纹库内查找具有一致时间偏移的匹配\cite{audfprint}。相比原波形逐点比较，地标结构能够利用录音中的局部稳定关系，并支持从歌曲片段检索整曲。系统对匹配哈希数、时间跨度和首选/次选优势设置门槛，依据不足时拒绝给出确定答案。

在线识曲将硬件采集结果接入更大范围的第三方音乐识别目录。项目完成了录音获取、片段生成、进程管理、节流、结果解析与本地回退，Shazam服务承担歌曲目录及其识别能力\cite{shazamio}。同一条任务记录同时保留所选模式与实际结果来源，避免将本地命中统计为在线命中。

\fig{intelligence}{从物理数据到语义与动作的智能感知闭环}{深色链路表示已实现能力；右侧给出可扩展AI方向，属于后续设计而非本次性能结论。}

从AI系统工程角度，自研贡献集中在FPGA感知前端、可靠数据、自然交互与应用调度，形成特征到决策反馈的闭环。第三方组件保留来源；学习式融合与边缘推理可复用数据接口，另行训练和独立评价。

\section{共享服务、显示隔离与体感实时性}
多应用系统的核心难点之一是设备所有权。若每个页面独立打开串口，切换时可能出现争用、重复命令或配置遗留。最终设计只保留一个串口拥有者；游戏使用消息订阅，识曲使用同一服务的任务控制。进入和退出游戏保存并核对触发配置，构成可逆的资源占用过程。

PANEL显示命令固定16字节，带会话、显示序号与录音号。旧会话、旧记录和较小序号不能覆盖新显示；相同序号只刷新心跳，不重复播放收藏动画。电脑心跳超过3秒失效后，主机相关阶段转入离线提示，板端真实录音计数及本地监测仍可继续。显示失败有独立超时和降级，不应导致完整录音读取失败。

体感游戏分别处理输入新鲜度、120 Hz逻辑与绘制；失联或离页时暂停，返回后手动继续。延迟仅统计主机消息发布至Canvas提交，明确软件测量边界。
